\section{Implementation}
\label{sec:implementation}

\subsection{Technology Stack}

The prototype is implemented primarily in Python and uses Apache Spark as the
structured data-processing engine.

The main technologies are:

\begin{itemize}

    \item \textbf{Python 3.12} for orchestration, catalog management,
    retrieval, validation, experimental evaluation, and reporting;

    \item \textbf{PySpark 3.5.8} and \textbf{Java 17} for Spark SQL execution
    over the curated Parquet data lake;

    \item \textbf{SQLite} for the persistent structured metadata catalog;

    \item \textbf{Qdrant} in local embedded mode for vector-based metadata
    retrieval;

    \item \textbf{Ollama} for running both the embedding model and the local
    Text-to-SQL language model;

    \item \texttt{embeddinggemma} for metadata embeddings;

    \item \texttt{qwen2.5-coder:3b} for Spark SQL generation and optional
    one-shot repair;

    \item \textbf{SQLGlot} for SQL parsing and dialect-aware validation;

    \item \textbf{Matplotlib} for generation of reproducible report figures.

\end{itemize}

All language-model inference is executed locally. No paid external model API
is required by the final prototype.


\subsection{Execution Environment}

The implementation and experiments were executed in a GitHub Codespaces
environment based on Ubuntu Linux.

Spark runs locally using two worker threads. The principal experimental
configuration uses:

\begin{verbatim}
--master local[2]
--driver-memory 1g
\end{verbatim}

Spark SQL uses four shuffle partitions in the controlled data-volume
benchmark, while adaptive query execution is enabled.

The local execution environment is intentionally modest. It contains two CPU
cores and no dedicated GPU. This constraint is relevant when interpreting
language-model latency results: the measured inference times describe the
CPU-only experimental environment and are not intended as general performance
claims for the underlying model.


\subsection{Repository Organization}

The implementation is organized into separate directories for configuration,
source code, scripts, datasets, and experimental artifacts.

The principal structure is:

\begin{verbatim}
config/
data/
  raw/
  curated/
  catalog/
artifacts/
  benchmarks/
  report/
scripts/
  catalog/
  data/
  experiments/
  reporting/
  retrieval/
src/
  metadata_agent/
report/
\end{verbatim}

The \texttt{src/metadata\_agent} package contains reusable application
components such as the catalog access layer, metadata document construction,
embedding support, relation-aware retrieval, SQL validation, and repair
logic.

The \texttt{scripts/} hierarchy contains reproducible entry points for data
preparation, metadata indexing, validation, experimental evaluation, and
final report artifact generation.

Configuration files are stored separately under \texttt{config/}. This keeps
experimental inputs, such as benchmark questions and synthetic catalog-scale
distractors, outside the implementation code.


\subsection{Curated Data Storage}

Raw source data and curated analytical datasets are stored separately.

The curated Spark datasets are written in Parquet format. This choice provides
a columnar representation directly supported by Spark SQL and avoids requiring
a separate database server for physical query execution.

The main curated paths correspond to January 2024 Yellow Taxi, Green Taxi,
Taxi Zones, and hourly weather data. These physical paths are recorded in the
metadata catalog together with schema, row-count, temporal, and semantic
information.

Generated raw and curated datasets are excluded from Git version control
because they are reproducible and potentially large. Experimental benchmark
summaries and final report artifacts, in contrast, are committed to the
repository.


\subsection{Metadata Catalog Implementation}

The metadata catalog is materialized as a SQLite database. Its content is
constructed from a versioned configuration file and the actual curated Spark
schemas.

The core tables include:

\begin{itemize}

    \item \texttt{datasets};
    \item \texttt{columns};
    \item \texttt{semantic\_concepts};
    \item \texttt{column\_semantics};
    \item \texttt{relationships};
    \item \texttt{aliases};
    \item \texttt{semantic\_rules};
    \item \texttt{catalog\_meta}.

\end{itemize}

Physical column definitions are checked against the curated Parquet schemas
during catalog construction. This reduces the risk of maintaining metadata
that refers to schema elements that do not exist in the execution layer.

The generated SQLite file is treated as a reproducible artifact rather than a
source file and is therefore excluded from Git.


\subsection{Metadata Document Construction and Embedding}

The vector index does not embed the complete catalog as a single document.
Instead, individual metadata entities are converted into separate textual
documents.

One document is generated for each physical column, semantic concept, dataset,
relationship, and semantic rule. The base configuration therefore produces
103 independently retrievable documents.

Each document includes enough natural-language description to support
semantic matching while its payload preserves the corresponding metadata
entity type and key.

Embeddings are produced locally using \texttt{embeddinggemma}. The resulting
768-dimensional vectors are inserted into a Qdrant collection using cosine
distance.

The index can be rebuilt deterministically from the metadata catalog and is
therefore not committed as a binary repository artifact.


\subsection{Relation-Aware Retrieval Implementation}

The relation-aware retriever is implemented as a deterministic Python
component.

For each question it combines:

\begin{enumerate}

    \item dense Top-\(k\) retrieval from Qdrant, with \(k=5\);

    \item lexical matching against metadata names and aliases;

    \item semantic-concept expansion toward physical columns;

    \item dataset selection using explicit and inferred evidence;

    \item relationship selection and endpoint expansion;

    \item semantic-rule expansion when supported by the selected metadata.

\end{enumerate}

Explicit dataset mentions receive stronger priority than unrelated dense
retrieval candidates. This prevents a semantically similar but physically
different dataset from replacing a dataset named directly by the user.

Relationship selection also uses query evidence to distinguish alternatives.
For example, pickup-oriented and drop-off-oriented relationships connect the
same taxi source with the same Taxi Zones source but require different
physical columns.

The retrieval implementation returns both a structured metadata selection and
a rendered SQL-safe context. This makes it possible to evaluate retrieval
quality independently from downstream LLM generation.


\subsection{SQL-Safe Context Rendering}

The renderer converts the internal metadata selection into the text included
in the LLM prompt.

The implementation explicitly distinguishes between metadata identifiers and
physical Spark SQL identifiers. Internal concept names, alias names,
relationship identifiers, and rule identifiers are not presented as if they
were executable columns or tables.

Physical schema objects are rendered using fully qualified names whenever
appropriate. Join metadata is rendered as executable physical conditions, and
semantic rules are rendered as executable predicates or expressions.

The renderer also adds constraints explaining that the registered January
2024 tables are already temporally restricted. This avoids encouraging the
model to invent an unnecessary year or month column when no narrower date
filter is requested.


\subsection{Local LLM Configuration}

Text-to-Spark-SQL generation is performed through Ollama using
\texttt{qwen2.5-coder:3b}.

The generation parameters are fixed to:

\begin{itemize}

    \item temperature \(=0\);
    \item context size \(=4096\) tokens.

\end{itemize}

The deterministic temperature setting reduces sampling variability between
benchmark runs.

The same model is used for the complete-catalog baseline and for the
relation-aware configurations. As a result, the main A/B comparison changes
the metadata context rather than the language model.

The repair configuration also uses the same local model, but its prompt
contains the failed SQL and validator feedback in addition to the original
question and retrieved context.


\subsection{SQL Validation Implementation}

Validation is implemented using both SQLGlot and Spark.

SQLGlot first parses the generated query using the Spark dialect. The
validator then checks that the request contains a single read-only query and
that referenced tables belong to the allowed physical dataset set.

Spark analysis provides the second validation layer. Because Spark resolves
the query against the registered schemas, it can identify physical problems
such as unknown columns, unknown relations, incompatible expressions, and
other analysis errors.

The benchmark also associates each question with an expected output schema.
This allows the validator to detect cases in which the generated query returns
the wrong result structure even if its SQL is otherwise executable.


\subsection{One-Shot Repair Implementation}

Repair is implemented as a separate component rather than as an unrestricted
conversation loop.

A repair is attempted only after validation failure. The repair prompt
contains:

\begin{itemize}

    \item the original natural-language question;

    \item the metadata context used for the original generation;

    \item the invalid SQL query;

    \item the validation error returned by the validator.

\end{itemize}

Only one repaired SQL candidate is generated. The candidate is passed through
the same validation stage before any execution is attempted.

This implementation deliberately prevents repair from using the gold query or
gold result. It therefore tests whether validator feedback and already
available metadata are sufficient to recover from a generation error.


\subsection{Benchmark Configuration}

Benchmark questions and gold specifications are stored as versioned JSON
configuration files.

Each benchmark item can include:

\begin{itemize}

    \item the natural-language question;

    \item required datasets;

    \item required physical columns;

    \item required relationships;

    \item required semantic rules;

    \item expected output columns;

    \item gold Spark SQL;

    \item gold result information.

\end{itemize}

This representation supports evaluation at multiple levels. Metadata
retrieval can be compared with the required metadata set, while generated SQL
can be validated and ultimately compared by result rather than only by exact
query text.


\subsection{Experimental Artifact Management}

Experimental outputs are written under \texttt{artifacts/benchmarks} as JSON
or text files. These files contain the final numerical results used by the
report.

Temporary execution logs and rebuildable indexes are stored separately and
excluded from version control.

The final reporting pipeline reads only the frozen benchmark JSON files. It
generates six figures in both PDF and PNG format and four summary tables in
both LaTeX and CSV format. It does not rerun Spark, Ollama, retrieval, or any
experimental benchmark.

This separation ensures that report generation cannot modify the underlying
experimental measurements.


\subsection{Reproducibility Strategy}

Reproducibility is supported at three levels.

First, the environment and Python dependencies are explicitly versioned,
including PySpark, SQLGlot, Qdrant client libraries, and Matplotlib.

Second, benchmark configurations, experiment scripts, and generated summary
results are committed to Git.

Third, the project uses Git tags to separate experimental protocol freezes
from observed results. Important evaluation stages were therefore committed
and tagged before the corresponding benchmark was executed.

This strategy is discussed in more detail in
Appendix~\ref{sec:reproducibility}.
